% !PRINTED 538
\weberpage{538}

Sind die $y$ dann so bestimmt, so setzen wir
\weq{10}{\omega = \omega_1 y_1 + \omega_2 y_2 + \cdots + \omega_n y_n,}
und leiten aus (5) die Gleichung ab
\weq{11}{\lambda\,\omega = \sum^{s} \alpha_s u_s.}

Da nun die $\alpha_s$ durch $\lambda$ und die $u_s$ durch $p$ theilbar sind,
so folgt, dass $\omega$ durch $p$ theilbar ist, und dass mithin, nach
\S.\,145, 2., $y_1, y_2, \ldots, y_n$ durch $p$ theilbar sein müssen, was
unserer Voraussetzung entgegen ist.

Damit ist bewiesen, dass $T$ eine primitive Function ist, \gesperrt{und
dass also der absolute Werth der Determinante $A$
gleich der absoluten Norm von $\lambda$ und folglich auch
von $\mu$ ist} (\S.\,138).

Wenden wir das Ergebniss auf die specielle Basis (3), \S.\,146
an, so ergiebt sich die Formel:
\weq{12}{N_a(\mu) = a_{1,1}\,a_{2,2}\ldots a_{n,n}.}

Hieran knüpfen wir noch folgende Bemerkungen: Wenn
man in einer Basisform eines Functionales $\mu$
\[
\lambda = \alpha_1 t_1 + \alpha_2 t_2 + \cdots + \alpha_n t_n
\]
auf die Variablen $t_1, t_2, \ldots, t_n$ eine ganzzahlige lineare Sub\-
stitution mit der Determinante $\pm 1$ anwendet, so entsteht eine
neue Basisform von $\mu$. Denn allen ganzzahligen rationalen Werthen
der Variablen $t_1, t_2, \ldots, t_n$ entsprechen ganzzahlige rationale
Werthe der neuen Variablen und umgekehrt.

Die Anwendung einer linearen Substitution auf die Variablen $t$
ist aber gleichbedeutend mit der Anwendung der transponirten
Substitution auf die Co\"efficienten $\alpha_1, \alpha_2, \ldots, \alpha_n$ (\S.\,37), d.\,h. mit
dem Uebergange zu einer neuen Basis von $\mu$:
\weq{13}{(\beta_1, \beta_2, \ldots, \beta_n) = (B)\,(\alpha_1, \alpha_2, \ldots, \alpha_n),}
die dann durch Zusammensetzung mit (1) ergiebt:
\weq{14}{(\beta_1, \beta_2, \ldots, \beta_n) = (B)\,(A)\,(\omega_1, \omega_2, \ldots, \omega_n).}

Wenn die Discriminante (\S.\,144) der durch eine Substitu\-
tion (13) bestimmten Zahlen $\beta$ mit der Discriminante der $\alpha$
übereinstimmt, so muss die Substitutionsdeterminante $B = \pm 1$
sein; und wir kommen also zu den Sätzen:
\begin{enumerate}
\item[1.] \gesperrt{Die Discriminante einer Basis von $\mu$ ist gleich
dem Quadrate der absoluten Norm von $\mu$, multi\-
plicirt mit der Grundzahl des Körpers;}
\end{enumerate}
%CONT>>
